\documentclass{report}
\usepackage{amsmath,amssymb}
\setlength{\parindent}{0mm}
\setlength{\parskip}{1em}
\begin{document}
\begin{center}
\rule{6in}{1pt} \
{\large Eloy Romero Alcalde \\
{\bf Spectrum Slicing in the Context of Computing Many Eigenvalues in Large Hermitian Matrices}}

Computer Science Department \\ College of William \& Mary \\ McGlothlin-Street Hall 108 \\ 251 Jamestown Rd \\ Williamsburg \\ VA 23185
\\
{\tt eloy@cs.wm.edu}\\
Andreas Stathopoulos\end{center}

We address the challenge of computing many eigenvalues and eigenvectors
of sparse Hermitian matrices so large that neither direct methods nor
unrestarted iterative methods such as Lanczos are feasible in terms of
memory space. Typically, problems with large dimension are solved with
restarted methods such as Thick-Restart Lanczos, Locally Optimal Block
Preconditioned Conjugate Gradient (LOBPCG), Generalized Davidson with CG
restarting (GD+k) and Jacobi-Davidson, which are competitive under
limited memory.

However the performance of the restarted methods does not scale linearly
with respect to the number of eigenvalues sought because of the
computational cost of keeping the new solutions orthogonal to the already
converged eigenvectors. As a result, orthogonalization introduces a
quadratic term with respect to the number of computed eigenpairs.

We consider spectrum slicing which attempts to improve scalability by
splitting the target spectrum region into subintervals and computing all
the eigenvalues in each of them independently. The approach avoids
orthogonalization against the solutions computed in other subintervals
and at the same time introduces a new level of parallelism. However a
practical implementation requires 1) a strategy to split the target
region that balances the computation cost in each subinterval; for
instance splitting the region into subintervals with similar number of
eigenvalues; and 2) heuristics to dynamically tune the solver's
parameters based on the density of the eigenvalues in the subinterval.

In the last few years there has been renewed interest in this technique
encouraged by the emergence of high-performance and scalable solvers for
computing all the eigenvalues inside a region in the spectrum. These
solvers are based on Contour Integration (e.g., FEAST) and polynomial
filtering (e.g., FILTLAN).

In this talk, first we focus on the choice of solver for within
individual intervals. We compare FEAST using iterative methods for the
solution of the linear systems and Davidson-type methods available in
PRIMME using polynomial filtering. We experiment with various polynomial
degrees and various levels of accuracy for the solvers. Our results show
that very small polynomial degrees and low linear solver accuracies are
better for best performance. They also show that JDMQR or GD+k using
polynomial filters outperform FEAST with iterative linear solvers.
Second, we show our strategies to automatically tune the solver and
choose the subintervals for spectrum slicing. These aim to find the best
performance by fine tuning the degree of the polynomial to balance the
total number of matrix-vector products versus the cost of
orthogonalization.


\end{document}
